% !TeX root = ../../ctufit-thesis.tex
%---------------------------------------------------------------
% CHAPTER 1: History of data centers
\chapter{Historie datových center a nároky moderní infrastruktury}

Digitální svět je často vnímán jako nehmotná entita existující v cloudu, ale jeho fyzické základy v podobě datových center procházejí nejvýraznější transformací od dob prvních sálových počítačů. První éru této infrastruktury charakterizuje nízkonákladová škálovatelnost pomocí serverů z běžných počítačových komponent, kdežto v současnosti datová centra tvoří na míru vyrobené počítače v měřítku skladu (WSC)~\cite[s.~11]{ranganathan_twenty_2024}. Paradigma datového centra se tak přesunulo od skladu informací k superpočítači provádějícímu náročné výpočty.

Právě tento technologický skok, doprovázený přechodem od chlazení vzduchem ke kapalinovým systémům~\cite{bizo_silicon_2022}, definuje nové nároky na infrastrukturu i její okolí (v podobě environmentální stopy). Následující text analyzuje historický vývoj těchto staveb, technologické příčiny jejich rostoucí náročnosti a faktory, které v současnosti určují rozmístění datových center na globální mapě.

\section{Historie a architektura}
\label{sec:historie}

%%%%%%%%%%%%%%%%%%%%%%%%%%%%%%%%%%%%%%%%%%%%%%%%%%%%%%%%%%%%%%%%%%%%%%%%%%%%%
% SEKCE 1: HISTORIE A ARCHITEKTURA
% - Definice datového centra napříč historií
% - Uvedení klasifikačního rámce Tier od Uptime Institute
% - Definice WSC a důvody přechodu od tradičních DC
% - Popis moderního AI datového centra + AI factory]
%%%%%%%%%%%%%%%%%%%%%%%%%%%%%%%%%%%%%%%%%%%%%%%%%%%%%%%%%%%%%%%%%%%%%%%%%%%%%

Pochopení současného technologického šoku vyžaduje retrospektivní pohled na to, jak se infrastruktura vyvíjela od improvizovaných začátků až po vysoce optimalizované systémy. Historicky byla datová centra chápána jako zázemí pro ukládání dat a distribuované výpočty, nicméně s rozvojem cloudu a následným nástupem specializovaných AI továren se tato role zásadně proměnila. Z původních ubytoven pro servery se staly integrované výpočetní celky, které dnes fungují jako jediné monolitické superpočítače. \cite{barroso_data_2026, ranganathan_twenty_2024, barroso_brief_2021}

\subsection{Tradiční pojetí DC a standardy spolehlivosti}
\label{subsec:tradice}

Rihards Balodis a Inara Opmane~\cite[s.~180--181]{tatnall_reflections_2012} v roce 2012 definovali datové centrum jako specializované fyzické prostředí určené pro bezpečné umístění počítačových systémů a jejich komponent. Jeho nezbytné zázemí tvoří systémy pro distribuci a zálohování elektřiny, redundantní komunikační rozvody, klimatizace, systémy požární ochrany a prvky fyzického zabezpečení objektu.

Modernější definice datových center zahrnují i samotná výpočetní zařízení. Například podle odborného článku z roku 2020 v časopise Facilities~\cite{fadaeefath_abadi_data_2020} je DC definováno jako významná infrastruktura obsahující velké množství serverů a počítačů, která poskytuje internetové služby společnostem po celém světě.

Starší definice ukazuje, že historicky byla datová centra vnímána především jako \textit{budovy}, jež mají poskytovat vhodné podmínky pro provoz počítačových systémů. Novější definice vnímá datové centrum nejen jako budovu, ale jako soubor vybavení pro zpracování, ukládání a přeposílání informací. Oba texty se však shodují v tom, že datová centra vyžadují speciální subsystémy pro napájení, regulaci okolní teploty a vlhkosti a komunikační nástroje~\cite{tatnall_reflections_2012, fadaeefath_abadi_data_2020}.

Výše uvedené klíčové subsystémy ovlivňují výslednou dostupnost služeb pro koncové uživatele. Pro sjednocení těchto požadavků navrhla organizace Uptime Institute čtyři třídy datových center v rámci celosvětově uznávané klasifikace Tier~\cite{turner_iv_tier_2006}. Tato certifikace definuje úrovně, které se od sebe liší zejména mírou redundance komponent a konfigurací distribučních cest pro provoz IT vybavení.

Souhrn parametrů pro jednotlivé úrovně Tier uvádí tabulka~\ref{tab:uptime_tiers}, detailnější popis je obsažen v příloze~\ref{app:tiery}. Tato klasifikace tvoří vrchol éry, v níž byla spolehlivost definována fyzickou redundancí komponent. S nástupem virtualizace se však pozornost přesunula k softwarově definované infrastruktuře.

\begin{table}[b] \centering \setlength{\tabcolsep}{4pt}
\begin{tabular}{|l|c|c|c|c|}
\hline
\textbf{Parametr} & \textbf{Tier I} & \textbf{Tier II} & \textbf{Tier III} & \textbf{Tier IV} \\ \hline
Dostupnost & 99,671~\% & 99,741~\% & 99,982~\% & 99,995~\% \\ \hline
Redundance komponent & N & N+1 & N+1 & >2N \\ \hline
Distribuční cesty & 1 & 1 & 1 akt., 1 zálož. & 2 akt. \\ \hline
Souběžná udržovatelnost & Ne & Ne & Ano & Ano \\ \hline
Odolnost proti poruchám & Ne & Ne & Ne & Ano \\ \hline
\end{tabular}
\caption{Souhrn požadavků pro klasifikaci Tier~\cite[s.~13]{turner_iv_tier_2006}. N označuje kritickou kapacitu nezbytnou pro podporu aktuálního IT zatížení.}
\label{tab:uptime_tiers}
\end{table}

\subsection{Koncept integrovaného systému WSC}
\label{subsec:wsc}

Na přelomu tisíciletí se začalo formovat nové paradigma distribuovaných výpočtů. Namísto delegování zátěže na koncová zařízení (osobní počítače), požadavky na výkon a kapacitu byly přesunuty do centralizované infrastruktury~\cite{barroso_data_2026}. Tento koncept poskytování škálovatelného výpočetního výkonu, softwaru a úložiště prostřednictvím internetu se nazývá Cloud Computing~\cite{lewis_basics_2010}.

Právě Cloud Computing umožnil vývoj softwarových systémů s nároky, které řádově přesahovaly možnosti izolovaných uživatelských zařízení -- úzkým místem se tak stala infrastruktura samotného provozovatele softwaru. V reakci na tyto potřeby vznikly mohutnější a výkonnější výpočetní celky s vyšší kapacitou, které L. A. Barroso definuje jako Warehouse-Scale Computers (WSC)~\cite{barroso_brief_2021, barroso_data_2026}.

Výpočetní výkon a velikost úložiště však nejsou jediné charakteristiky, které WSC odlišují od klasických datových center. Barroso rozdíl ilustruje v pěti dimenzích hardwarové a softwarové architektury~\cite[s.~2--4]{barroso_data_2026}:

\begin{description}
    \item[Homogenita] Tradiční datová centra hostí různorodý hardware mnoha firem. WSC zpravidla využívá jednotný hardware pod kontrolou jediné organizace, aby byla zajištěna rychlá nahraditelnost a nákladová efektivita.
    
    \item[Software] Zatímco běžný poskytovatel dodává pouze budovu, konektivitu a fyzickou správu serverů~\cite[s.~10]{marx_gomez_foundations_2017}, WSC disponuje komplexní softwarovou vrstvou, která tisíce serverů propojuje a umožňuje jim komunikovat.
    
    \item[Škálovatelnost] WSC je navržen pro běh masivních internetových služeb využívajících desetitisíců serverů současně. To je umožněno díky efektivní komunikaci mezi servery, kterou běžná DC neposkytují.
    
    \item[Odolnost] Tradiční centra se snaží chybám hardwaru předcházet, což je nákladné. WSC naopak s výpadky komponent počítá a řeší je pomocí architektury navržené pro automatickou (graceful) toleranci jejich poruch.
    
    \item[Měřítko] WSC bývají o jeden až dva řády větší než běžná podniková datová centra, přičemž veškeré prostředky (dodávky energie, chlazení) jsou sdíleny v rámci společné infrastruktury pro správu zdrojů.
\end{description}

Technologický vývoj v tomto období (roky 2004 až 2008) se soustředil také na optimalizaci provozu. Důležitou inovací bylo zavedení metriky PUE, která je popsána v kapitole~\ref{subsec:pue}. S tím byl kladen důraz i na energetickou škálovatelnost -- schopnost hardwaru měnit svou spotřebu podle aktuálního zatížení.~\cite{ranganathan_twenty_2024}

Z hlediska budoucího vývoje je zásadní období mezi lety 2009 a 2013, kdy se začaly objevovat první úlohy zaměřené na hluboké učení (např. projekt Google Brain)~\cite{ranganathan_twenty_2024}. Tento moment předznamenal přerod univerzálních WSC v úzce specializované AI továrny, které dominují současnému trhu.

\subsection{Architektura moderních AI továren}
\label{subsec:ai_factory}

Nejmodernější druh datových center se nazývá AI Factory (továrna na AI). Podle slovníku společnosti NVIDIA, předního dodavatele hardwaru pro AI datová centra, jde o specializovanou výpočetní infrastrukturu, která vytváří hodnotu spravováním dat po celou dobu životního cyklu AI.~\cite{noauthor_nvidia_nodate}

Popisovaný životní cyklus je iterativní proces zahrnující sbírání dat, trénování a ladění modelu až po jeho hodnocení. Poté se model přesouvá do fáze inference, tedy přeměny vstupu na výstupy (dotazování modelu). Koncept AI továrny představuje moderní přístup, který se snaží tyto fáze propojit v jednom datovém centru. V praxi se však prozatím provozují zpravidla odděleně:

\begin{itemize}
    \item Trénovací centra jsou navržena pro vysoký a konzistentní výkon, pročež vyžadují pokročilé metody chlazení~\cite{cruzes_data_2025}. Jde o průlomovou technologii.

    \item Inferenční centra kladou důraz na efektivitu a nízkou latenci, ale zátěž je vysoce stochastická~\cite{cruzes_data_2025, ginzburg-ganz_technical_2025}. Podobají se WSC a jsou již dobře zavedená.
\end{itemize}

Rozdíl mezi trénovacími centry a WSC spočívá v přechodu od izolovaných úloh k realizaci masivně paralelizovaných výpočtů. Barrosova definice WSC počítá s nezávislostí jednotlivých serverů (jakýkoli server lze díky virtualizaci nahradit). Naproti tomu TC funguje jako jeden logický superpočítač, takže porucha jediného procesoru může zastavit trénování celého modelu~\cite[s.~270]{barroso_data_2026}. Tato provázanost definuje trénovací centrum jako úzce specializovaný stroj, kde je integrace a rychlost propojení důležitější než kdykoli předtím.

Dalším rozdílem je požadovaná dostupnost: Zatímco tradiční datová a inferenční centra bývají zpravidla Tier III nebo IV, trénovací centra mohou využívat architekturu s nižší úrovní fyzické redundance, často odpovídající úrovni Tier I. Případný výpadek pro ně díky postupnému ukládání stavu nepředstavuje kritické ohrožení služby, ale pouze dočasnou ztrátu výpočetního času~\cite[s.~83]{barroso_data_2026}. Tato optimalizace pro nákladovou efektivitu a rychlost výstavby na úkor drahé dostupnosti umožňuje zkrátit čas potřebný pro vývoj nových modelů, což je pro rychle se vyvíjející odvětví z byznysového hlediska důležité~\cite[s.~24]{ginzburg-ganz_technical_2025}.

\section{Konstrukce a hardware}
\label{sec:konstrukce}

%%%%%%%%%%%%%%%%%%%%%%%%%%%%%%%%%%%%%%%%%%%%%%%%%%%%%%%%%%%%%%%%
% SEKCE 2: KONSTRUKCE A HARDWARE
% - Vnitřnosti DC - racky, servery, procesory, switche, ...
% - Přechod ke specializovaným čipům, CPU -> GPU -> TPU
% - Uspořádání v datovém centru, co obvykle zahrnují z prostorového hlediska
%%%%%%%%%%%%%%%%%%%%%%%%%%%%%%%%%%%%%%%%%%%%%%%%%%%%%%%%%%%%%%%%

Tato část se zaměřuje na jednotlivé komponenty tvořící a podporující výpočetní výkon datových center. Stěžejním tématem je standardizace serverových skříní, fyzická konstrukce moderních akcelerátorů a členění vnitřních prostorů datového centra, které musí splňovat extrémní nároky na koncentraci techniky a statické zatížení budov. Zdrojům energie a chlazení DC se věnují sekce \ref{sec:energie} a \ref{sec:chlazeni}.

\subsection{Fyzické parametry serverových skříní}
\label{subsec:racky}

Základním stavebním kamenem každého datového centra je tzv. rack. Jde o standardizovaný kovový rám, do kterého se horizontálně zasouvají a upevňují modulární komponenty (servery, úložiště, síťové přepínače, směrovače, \dots{}). Vnější konstrukce standardního racku je 78 palců vysoká, 23--25 palců široká a 26--30 palců hluboká\footnote{Přibližně 1 980~mm vysoká, 580--640~mm široká a 660--760~mm hluboká.}, ale existují i jiné velikosti.~\cite[s.~2940]{kant_data_2009}

Pro praktické použití racků jsou nejdůležitější montážní šířka a výška. Podle standardu EIA-310-D má mít běžný rack montážní šířku 19 palců (482,6~mm), což odpovídá standardní šířce předního panelu IT vybavení. Stejný standard definuje jednotku U (1~U = 44,45~mm), která představuje montážní výšku typického IT zařízení. Pokud je tedy rack popsán jako 42U, dokáže pojmout 42 zařízení s výškou 1~U. Specifická zařízení, zejména výkonné AI akcelerátory, však mohou mít kvůli chlazení výšku 4~U, 8~U nebo i více.~\cite{hu_how_2020}

Z těchto důvodů standardní racky pro potřeby výkonných AI serverů nepostačují. V. Avelar~\cite[s.~13--14]{avelar_ai_2023} doporučuje, aby byly racky pro trénování umělé inteligence širší (např. 750~mm místo 600~mm pro uložení rozvodů kapalinového chlazení), hlubší (přes 1200~mm) a vyšší (48~U a více). Plně osazený AI rack může vážit přes 900~kg, což je potřeba zvážit při jeho přepravě i návrhu DC.

\subsection{Čipy a specializované akcelerátory}
\label{subsec:hardware}

Historicky byla datová centra poháněna univerzálními procesory (CPU), které jsou vynikající pro sériové zpracování složitého, větveného a nepravidelného kódu. Trénování umělé inteligence je však výpočetně velmi specifické -- spoléhá na několik málo operací, které jsou paralelně aplikovány na obrovské množství dat~\cite[s.~133--134]{barroso_data_2026}. CPU tak byly nahrazeny jinými druhy čipů, které jsou pro takto specifické úkoly efektivnější~\cite[s.~1391]{phani_suresh_paladugu_evolution_2025}.

Mezi používané komponenty patří především grafické procesory (GPU), které obsahují tisíce jednodušších jader (architektura SIMT) užitečných při paralelizaci. Dalším evolučním krokem jsou čipy navržené čistě pro potřeby AI, což je činí efektivnějšími. Za zmínku stojí TPU (Tensor Processing Units) vyvinuté společností Google, které využívají systolická pole -- propojenou síť výpočetních jednotek optimalizovaných pro maticové násobení.~\cite[s.~134--139]{barroso_data_2026}

Některé analogie uvádějí, že CPU je jako sportovní auto, protože dokáže velice rychle přepravit pár lidí. Naproti tomu GPU jsou jako autobusy, které sice jedou pomaleji a nevejdou se na všechny cesty, ale dokážou přepravit stovky lidí najednou. Nejsou kvalitativně horší ani lepší, jen slouží k jinému účelu.

\subsection{Vnitřní uspořádání a technické zóny}
\label{subsec:dispozice}

Vnitřní dispozice datového centra se člení na dvě hlavní oblasti: White Space a Gray Space. První představuje plochu vyhrazenou pro IT vybavení (racky, síťové prvky a úložná pole). Gray Space naopak zahrnuje plochu pro veškerou podpůrnou infrastrukturu (transformátory, UPS a systémy chlazení)~\cite[s.~300]{geng_data_2014}. U tradičních DC tvořilo technické zázemí menší část celkové dispozice, přičemž poměr mezi Gray a White Space byl u větších datových center až 1:3~\cite{butler_unlocking_2025}.

Extrémní výkonová hustota AI racků však tento poměr mění: Jednak zmenšuje nároky na podlahovou plochu v rámci počítačového sálu (výkon je mnohem více koncentrovaný), jednak vyžaduje větší napájecí a chladicí systémy, které fyzicky zabírají stále více prostoru (zejména kvůli mohutnému potrubí). Jejich vzájemný poměr se tak může posunout až nad hranici 1:1~\cite{butler_unlocking_2025}.

\section{Energie}
\label{sec:energie}

%%%%%%%%%%%%%%%%%%%%%%%%%%%%%%%%%%%%%%%%%%%%%%%%%%%%%%%%%%%%%%%%
% SEKCE 3: ENERGIE
% - Co je to příkon a jaký je typický příkon jednoho datového centra
% - Jak se měří efektivita DC -- PUE
% - Odkud datová centra berou energii (důležité pro další kapitoly)
% - TODO: Možná přidat sekci o dynamické spotřebě AI továren - z minuty na minutu přestanou odebírat megawatty energie a síť nestíhá reagovat. To by ale asi bylo dlouhé prodloužení, protože je aspoň trochu potřeba vysvětlit, jak funguje energetická síť. 
%       - Pozor: Napsal jsem to do shrnutí, takže pokud se rozhodnu že nebudu psat tento odstavec, musím to ze shrnutí odstranit.
%%%%%%%%%%%%%%%%%%%%%%%%%%%%%%%%%%%%%%%%%%%%%%%%%%%%%%%%%%%%%%%%

Energetická infrastruktura se stala hlavním limitujícím faktorem moderních datových center. Tato kapitola analyzuje, jakým způsobem revoluce v oblasti umělé inteligence narušila dekádu trvající stabilitu energetické spotřeby a jaké nároky klade nejen na samotná DC, ale také na sdílenou energetickou síť.

\subsection{Příkon a spotřeba datových center}
\label{subsec:prikon}

Při analýze energetické náročnosti je nezbytné rozlišovat mezi příkonem (jednotka watt) a spotřebou (jednotka watthodina). V tomto kontextu představuje příkon okamžitý odběr elektřiny, který v čase kolísá podle vytížení DC. Instalovaný příkon definuje horní hranici odběru při 100\% vytížení. Spotřeba vyjadřuje množství elektřiny, která byla odebrána za určitý časový úsek.

Zatímco menší podniková datová centra měla historicky příkon menší než 1~MW~\cite{barroso_data_2026, sheng_power_2026}, větší datová centra určená pro Cloud Computing mají příkon 10 až 50~MW~\cite[s.~38]{iea_energy_2025}. Moderní hyperscale AI trénovací centra však fungují v úplně jiných mezích -- jedno takové zařízení vyžaduje příkon 100~MW až 1 GW~\cite{chen_electricity_2025, sheng_power_2026}. Přehled instalovaného příkonu datových center je uveden v tabulce~\ref{tab:energeticke_rozsahy}.

\begin{table}[t]
    \centering
    \begin{tabular}{|l|c|l|}
        \hline
        \textbf{Typ DC} & \textbf{Instal. příkon} & \textbf{Zaměření} \\ \hline
        Enterprise DC & < 1~MW & Lokální firemní IT \\ \hline
        Cloud Computing & 10--50~MW & SaaS, IaaS \\ \hline
        Hyperscale AI & 100--1~000+~MW & Trénování LLM \\ \hline
    \end{tabular}
    \caption[Srovnání instalovaného příkonu]{Srovnání instalovaného příkonu dle typu datového centra. Vlastní zpracování podle dat od Barrosa~\cite{barroso_data_2026}, IEA~\cite{iea_energy_2025} a Chen~\cite{chen_electricity_2025}.}
    \label{tab:energeticke_rozsahy}
\end{table}

Je nutné zvážit skutečné využití instalovaného příkonu. Zatímco u cloudových DC je typické využití asi 33,5~\%~\cite[s.~224]{barroso_data_2026}, u AI trénovacích center se blíží 100~\%. Jinými slovy, po dobu trénování modelu (které může trvat až měsíce) se průměrný příkon datového centra blíží k jeho instalovanému příkonu~\cite[s.~3]{avelar_ai_2023}. Příkladem: Pokud ve 200MW datovém centru běží AI servery 80~\% času~\cite[s.~27]{LBNL-2001637} (zbylých 20~\% jsou prostoje), jeho roční spotřeba činí asi 1,401~TWh. To odpovídá spotřebě Karlovarského kraje za rok 2024 ($\approx$1,5~TWH)~\cite{energeticky_regulacni_urad_spotreba_2024}.

\subsection{Power Usage Effectiveness}
\label{subsec:pue}

Pro zvýšení efektivity je nutné co největší část spotřebované energie využít na samotný výpočetní výkon, tedy na reálné zpracování dat a provoz výpočetních uzlů. V ideálním případě by měla většina energie směřovat přímo k serverům, akcelerátorům a dalším prvkům, které se podílejí na trénování modelu či poskytování služeb uživatelům. V praxi však tohoto stavu nelze dosáhnout bez odpovídající podpůrné infrastruktury, která zajišťuje napájení, chlazení, fyzickou bezpečnost i základní provozní stabilitu celého datového centra.

Z tohoto důvodu je nutné podrobně analyzovat, kam všude energie v datových centrech skutečně proudí, a jaký podíl připadá na jednotlivé části systému. Spotřebu energie v datových centrech proto dále rozdělujeme na~\cite{chen_electricity_2025, sheng_power_2026}:

\begin{itemize}
    \item \textbf{IT vybavení} (servery, disková pole, sítě): zhruba 40--50~\% spotřeby.

    \item \textbf{Chladicí infrastruktura}: tradičně spotřebovává 30--40~\% dodané energie.

    \item \textbf{Podpůrná zařízení} (napájení, osvětlení, zabezpečení): zbylých 10--30~\%.
\end{itemize}
 
K měření celkové efektivity se používá metrika PUE (Power Usage Effectiveness), která udává poměr mezi celkovou energií dodanou do budovy a energií, která skutečně dorazí k IT vybavení~\cite{chen_electricity_2025}:

$$\text{PUE} = \frac{\text{Total Facility Electricity Consumption}}{\text{IT Equipment Electricity Consumption}} \geq 1.$$

Zatímco v roce 2010 byla průměrná hodnota PUE datových center po celém světe kolem 2,5 (tzn. na každý 1 watt pro server připadlo 1,5 wattu na chlazení a ztráty), dnes se ji podařilo snížit na 1,55~\cite{de_roucy-rochegonde_ai_2025}. Moderní AI továrny však míří k vyšší efektivitě -- Google a Meta uvádějí, že jejich vybraná DC mají PUE pod 1,1 ~\cite{chen_electricity_2025, sheng_power_2026}, pravděpodobně díky efektivnímu chlazení vodou.

\subsection{Odkud datová centra berou energii}
\label{subsec:odkud_energie}

%%%%%%% STOPPED PROOFREADING HERE %%%%%%%

Většina datových center je primárně napájena z veřejné elektrické sítě. Hyperscale centra se často připojují přímo k vysokonapěťové přenosové soustavě (v USA 115--230 kV), zatímco menší zařízení využívají sítě středního napětí~\cite[s.~4]{chen_electricity_2025}. Původ této energie tak plně závisí na energetickém mixu daného regionu, přičemž celosvětově v napájení DC převažují fosilní paliva. Podíl obnovitelných zdrojů však má v budoucnu růst~\cite[s.~86]{iea_energy_2025}.

Z důvodu přetížení veřejných sítí, zdlouhavých povolovacích procesů a snahy o vyšší spolehlivost však provozovatelé stále častěji přistupují k vlastní výrobě a budování tzv. mikrosítí. Datová centra mohou integrovat lokální obnovitelné zdroje, jako jsou střešní solární panely nebo instalace větrných turbín, aby snížila svou závislost na centrální síti.~\cite[91--93]{barroso_data_2026}

Kromě toho je každé centrum vybaveno vlastními záložními generátory (historicky naftovými, nověji i plynovými) a v provozu se objevují i velkokapacitní palivové články (např. na vodík nebo bioplyn), které mohou sloužit nejen jako záloha, ale i jako primární zdroj energie. ~\cite{chen_electricity_2025}

Obrovská a nepřetržitá spotřeba umělé inteligence navíc podněcuje nový trend umisťování center přímo ke zdrojům: velcí hráči začínají stavět datová centra v bezprostřední blízkosti stávajících jaderných elektráren, nebo plánují nasazení malých modulárních reaktorů (SMR) přímo ve svých kampusech, aby získali stabilní bezuhlíkový výkon, aniž by zatěžovali či limitovali veřejnou síť.~\cite{chen_electricity_2025, cruzes_data_2025, iea_energy_2025} 

\section{Voda a chlazení}
\label{sec:chlazeni}

%%%%%%%%%%%%%%%%%%%%%%%%%%%%%%%%%%%%%%%%%%%%%%%%%%%%%%%%%%%%%%%%
% SEKCE 4: CHLAZENÍ
% - TDP jakožto ukazatel zvyšujících se nároků na chlazení
% - Nutnost přechodu od vzduchového ke kapalinovému chlazení
% - Dopady na spotřebu a odběr vody DC
% - TODO: Snažit se zakomponovat dodatky do normálních kapitol
%%%%%%%%%%%%%%%%%%%%%%%%%%%%%%%%%%%%%%%%%%%%%%%%%%%%%%%%%%%%%%%%

Prudký nárůst výkonu způsobil, že chlazení vzduchem přestalo kvůli jeho fyzikálním vlastnostem stačit. Kvůli tomu byli designéři datových center v podstatě nuceni přejít na kapalinové chlazení, protože kapaliny disponují mnohem vyšší tepelnou kapacitou a vodivostí. Dokáže tak odvést více tepla za kratší čas. To však způsobilo nový problém -- datová centra se stala závislá na dalším sdíleném zdroji, který je potřeba odněkud čerpat. 

\subsection{Thermal Design Power}
\label{sec:tdp}

S rostoucím výkonem čipů narůstá také množství generovaného tepla, které je nutné efektivně odvádět. K tomu slouží metrika TDP (Thermal Design Power) definující maximální tepelný výkon jednoho čipu, jenž musí chladicí systém kontinuálně odvádět při běžném provozním zatížení~\cite{mahajan_cooling_2006}. Hodnota TDP přibližně odpovídá doporučené horní hranici příkonu čipu~\cite{avelar_ai_2023}, což je v souladu s prvním termodynamickým zákonem. Nejedná se však o absolutní maximum -- při intenzivních výpočtech může reálný odběr hodnotu TDP krátkodobě překračovat, stejně jako v případě systematického přetaktování~\cite{noauthor_tdp_nodate}.

Nástup specializovaných AI čipů způsobil rychlý nárůst TDP u IT komponent, který musí moderní datová centra brát v potaz. Nejvýkonnější čipy dnes dosahují TDP až 1400~W. Pro srovnání, výkonné servery s CPU měly okolo roku 2010 TDP kolem 90 až 120~W, moderní mají asi 400~W~\cite[s.~3]{bizo_silicon_2022}. Přehled TDP několika vybraných modelů GPU je v tabulce~\ref{tab:gpu_tdp}.

\begin{table}[]
    \centering
    \begin{tabular}{|l|c|c|c|} \hline
        \textbf{GPU} & \textbf{Paměť} & \textbf{Rok vydání} & \textbf{TDP} \\ \hline
         V100 SXM2 & 32~GB & 2018 & 300~W \\ \hline
         A100 SXM & 80~GB & 2020 & 400~W \\ \hline
         H100 SXM & 80~GB & 2022 & 700~W \\ \hline
         B200 (Blackwell) & 186~GB & 2025 & 1000~W \\ \hline
         B300 (Blackwell Ultra) & 279~GB & 2025 & 1400~W \\ \hline
    \end{tabular}
    \caption[Přehled vývoje TDP]{Přehled vývoje TDP grafických procesorů NVIDIA. Tabulka převzata od Schneider Electric~\cite{avelar_ai_2023} a doplněna o nové údaje z dokumentace NVIDIA~\cite{noauthor_nvidia_2025, noauthor_nvidia_2025-1}.}
    \label{tab:gpu_tdp}
\end{table}

Extrémní TDP samotných čipů však není tím největším problémem. Kvůli požadavku na rychlou synchronizaci musí být tyto čipy umístěny fyzicky bezprostředně vedle sebe, jinak by bylo nutné použít velmi drahá~\cite[s.~4]{avelar_ai_2023} a energeticky náročná~\cite[s.~36]{cruzes_data_2025} optická vlákna. Z těchto důvodů se návrháři snaží umístit co nejvíce čipů do jednoho fyzického bloku. Instalují se tak do těsných formátů (například 8 GPU s vnitřní superrychlou sběrnicí NVLink v jednom serveru). 

Tato nevyhnutelná fyzická blízkost vede k prostorové koncentraci tepla, kvůli čemuž AI racky s nahuštěnými akcelerátory dosahují TDP 30 až 80~kW, u nejnovějších architektur i přes 100~kW  Taková hustota v jednom standardizovaném rámu absolutně přesahuje fyzikální možnosti vzduchového chlazení a vynucuje si přechod na chlazení kapalinové.~\cite[s.~23]{cruzes_data_2025}

\subsection{Od vzduchového chlazení ke kapalinovému}
\label{subsec:vzduch}

Tradiční vzduchové chlazení historicky dominovalo díky své jednoduchosti a nízké ceně. Spolehlivě funguje pro běžné serverové racky s hustotou výkonu zhruba do 20~kW~\cite[s.~10]{avelar_ai_2023}. S příchodem AI akcelerátorů, kde spotřeba jednoho racku stoupá na 30 až 100~kW, však vzduch narazil na své fyzikální limity a kapalinové chlazení se stalo absolutní nutností, protože odvádí teplo efektivněji než vzduch.~\cite[s.~4]{chen_electricity_2025} To přináší několik faktorů, které je při návrhu datových center potřeba zvážit:

\begin{description}
    \item[Výhody vodního chlazení] Systémy chlazení, kde je kapalina přiváděna pří\-mo k čipu (tzv. Direct-to-Chip neboli D2C architektura), umožňují snížit spotřebu energie na chlazení až o 90~\%. Kapalinové chlazení je tedy velmi efektivní - jak již bylo zmíněno, nejnovější datová centra používající tuto technologii dosahují PUE nižší než 1,1~\cite{chen_electricity_2025, sheng_power_2026}. Nejnovější systémy, jako je immersion cooling (servery jsou plně ponořené v kapalině), dosahují ještě vyšší efektivity (PUE < 1,03) a oproti vzduchovému chlazení mají o 95~\% nižší spotřebu energie.~\cite[s.~17--19]{cruzes_data_2025}
    \item[Nevýhody vodního chlazení] Vyšší účinnost je však vykoupena vysokými počátečními investicemi (CAPEX) a vyšší technologickou náročností. Pokročilé systémy jako D2C vyžadují specializované chladicí desky na procesorech, utěsněná potrubí uvnitř racků a jednotky pro distribuci chladivé kapaliny~\cite[s.~18]{cruzes_data_2025}. Navíc je velmi obtížné přestavět vzduchem chlazená DC na vodní chlazení, neboť často nemají dostatečně vyztužené podlahy (kapalina přináší dodatečnou zátěž)~\cite[s.~11]{avelar_ai_2023}.
\end{description}

Ačkoliv je tedy konstrukce vodního chlazení mnohem nákladnější, z dlouhodobého horizontu se tato investice vrátí v nižších nákladech na elektřinu. Především se ale jedná o technologickou nutnost, která je vynucena zvyšující se hustotou výkonu na jeden rack. Nabízí se ale otázka, kde datová centra berou vodu na chlazení a zda tento odběr nepřináší další náklady.

\subsection{Odběr a spotřeba vody moderních datových center}
\label{subsec:voda}

Přechod od vzduchového chlazení k pokročilým kapalinovým systémům dramaticky snížil spotřebu elektrické energie nutnou pro chlazení. Tento zisk v energetické efektivitě je však často vykoupen novým, velmi závažným problémem: zvýšenými nároky na spotřebu vody. Při hodnocení vlivu datových center na vodní zdroje je nutné rozlišovat mezi dvěma pojmy:

\begin{itemize}
    \item odběrem vody (water withdrawal), který představuje celkový objem sladké vody odčerpané z podzemních nebo povrchových zdrojů;
    \item a spotřebou vody (water consumption), která označuje nenávratně ztracenou část odebrané vody (nejčastěji se odpařuje do atmosféry);~\cite{li_making_2025}
    \item platí tedy vztah $\text{Water consumption} \le \text{Water withdrawal}.$
\end{itemize}

Datová centra odebírají vodu primárně pro své chladicí systémy. Teplo ze serverů se přenáší do vodního okruhu a v chladicí věži se část této vody cíleně odpaří, čímž se teplo uvolní do okolí. Protože se však při odpařování čisté vody v okruhu koncentrují minerály, soli a bakterie, musí se část zbylé vody pravidelně vypouštět a neustále doplňovat novou, čistou sladkou vodou. Z toho důvodu tyto systémy vyžadují neustálý přítok vody, často v kvalitě běžné pitné vody z vodovodu.~\cite{li_making_2025} Tomuto procesu se říká přímý odběr vody (Scope 1).

Mnohem větší zátěž na vodní zdroje často představuje tzv. nepřímý odběr (Scope 2). Jde o vodu odebíranou v elektrárnách, které pokrývají energetickou spotřebu datových center - zejména tepelné či jaderné elektrárny potřebují k výrobě elektřiny masivní chlazení a vodní elektrárny ztrácejí vodu odparem z nádrží~\cite{li_making_2025, LBNL-2001637}. V prostředí USA je průměrný nepřímý odběr vody pro výrobu elektřiny odhadován na 43,8 litrů na každou spotřebovanou kilowatthodinu~\cite{li_making_2025}.

Pro představu, typické 100MW hyperscale datové centrum může spotřebovat kolem 2 milionů litrů vody denně, což odpovídá spotřebě zhruba 6 500 domácností. 60~\% z této spotřeby je nepřímá.~\cite{li_making_2025} Vzhledem k tomu, že mnoho nových datových center vzniká v oblastech, které již dnes trpí nedostatkem vody (např. v Texasu nebo Arizoně), stává se tento masivní odběr sladké vody zdrojem konfliktů.

\subsection{Dodatky k odběru a spotřebě vody}
\label{subsec:dodatky_voda}

Dalším skrytým odběrem (Scope 3) je voda používaná v dodavatelském řetězci k výrobě samotného hardwaru. Ten je měřitelný jen velmi složitě, ale pravděpodobně je významný. Například společnost Apple uvádí, že 99~\% jeho vodní stopy pochází právě od dodavatelů hardwaru~\cite{li_making_2025}.

K měření efektivity hospodaření s vodou se dnes používá metrika WUE (Water Usage Effectiveness), která udává spotřebu vody v litrech na 1 kWh energie spotřebované samotným IT vybavením~\cite{jegham_how_2025}. Není však definována tak jasně jako její protějšek PUE, neboť v literatuře může být zahrnuta pouze spotřeba Scope 1, nebo i součet spotřeby přes Scope 1 až po Scope 3. Dále je možné namísto spotřeby použít odběr vody, kvůli čemuž je s touto metrikou obtížné pracovat a v této práci je uvedena pouze pro kontext. 

Celosvětový roční odběr vody spojený s provozem datových center se odhaduje na stovky miliard litrů a do roku 2030 by mohl přesáhnout 1 200 miliard litrů~\cite[s.~242]{iea_energy_2025}. Spotřeba na úrovní USA a celého světa bude dále analyzována v druhé kapitole.

\section{Lokalita}
\label{sec:lokalita}

%%%%%%%%%%%%%%%%%%%%%%%%%%%%%%%%%%%%%%%%%%%%%%%%%%%%%%%%%%%%%%%%%%%%%%%%%%%%%
% SEKCE 5: LOKALITA
% - Faktory ke zvážení při výstavbě datového centra
% - Formulováno jako PESTLE analýza pro lepší přehlednost, je to poměrně známý nástroj
%%%%%%%%%%%%%%%%%%%%%%%%%%%%%%%%%%%%%%%%%%%%%%%%%%%%%%%%%%%%%%%%%%%%%%%%%%%%%

Výběr lokality pro stavbu datového centra je komplexní proces, který funguje spíše jako vylučovací metoda než prosté hledání - začíná se širokým geografickým záběrem a postupně se eliminují místa, která nesplňují kritické požadavky, až zbyde několik málo optimálních kandidátů~\cite{geng_data_2014}. S nástupem AI datových center se navíc kritéria radikálně proměnila. Ty nejdůležitější byly autorem této práce rozděleny do skupin inspirovaných analýzou PESTLE:

\subsection{Politické faktory}
\label{subsec:pestle_po}

Při mezinárodní expanzi podniky zohledňují zákony o ochraně soukromí, cenzuru a geopolitická rizika spojená s tím, kde jsou citlivá data fyzicky uložena~\cite[s.~96]{geng_data_2014}. Důležité jsou také možnosti vlastnické struktury v daném státě, protože datová centra vyžadují obrovské investice.

Tyto faktory je zajímavé zvážit také z pohledu vlád: Protože se AI postupně stává záležitostí státní bezpečnosti, jednotlivé státy se mohou snažit snížit svou závislost na zahraničních datových centrech, což může způsobit jejich širší distribuci po celém světě~\cite[s.~25]{pilz_compute_2023}. 

\subsection{Ekonomické faktory}
\label{subsec:pestle_ec}

Náklady na výstavbu a provoz datového centra drasticky ovlivňují místní vlády, zejména na území USA. Provozovatelé proto hledají daňové pobídky a úlevy, například odpuštění daně z nemovitosti nebo daně z prodeje drahého IT vybavení.~\cite[s.~97--89]{geng_data_2014}

Moderní datová centra se stále častěji stavějí dále od velkých metropolí nejen kvůli jejich vysoké spotřebě energie, ale také možnosti postavit novou infrastrukturu, což je v zastavených oblastech téměř nemožné. Dalším faktorem je, že v odlehlých oblastech je jednodušší a rychlejší sehnat stavební povolení.~\cite[s.~25]{pilz_compute_2023}

Příkladem je Severní Virginie a její Data Center Alley, která má díky masivním úlevám a levné energii největší koncentraci datových center na světě. Dalšími významnými lokalitami jsou okolí Dallasu, Los Angeles, Chicaga nebo New Yorku.~\cite{pilz_compute_2023}

\subsection{Sociální faktory}
\label{subsec:pestle_so}

Přestože jsou datová centra silně automatizovaná, ke své stavbě a provozu vyžadují vysoce kvalifikovaný personál (stavbaře, IT inženýry, bezpečnostní techniky). Výhodou je blízkost velkých technických univerzit a výzkumných center~\cite{geng_data_2014}. Problém se stavaři v dnešní době nicméně vyřešen dočasnou relokací stávajících pracovníků stavebních firem, kteří cestují klidně přes celé Spojené státy na místo výstavby datového centra.

% TODO: Na tohle je potřeba najít lepší zdroje než youtube videa
Mnohem důležitějším aspektem jsou iniciativy obyvatel měst, ve kterých má být datové centrum postaveno -- výstavba se čím dál častěji setkává s nevolí obyvatel, kteří by žili v bezprostředním okolí datového centra. Lidé, kteří již v okolí DC žijí, si stěžují na neustálé hučení, které je způsobeno chladicími systémy. Datová centra mají také dopady na cenu elektřiny ve svém okolí, o čemž pojednává kapitola 2.

\subsection{Technologické faktory}
\label{subsec:pestle_te}

Dostupnost obrovského množství spolehlivé a levné elektrické energie je dnes hlavním limitem. Protože výstavba nových přenosových soustav trvá 4 až 10 let, datová centra narážejí na přetížené sítě a dlouhé pořadníky pro připojení (fronty mohou trvat i roky).~\cite{ginzburg-ganz_technical_2025, iea_energy_2025} Z tohoto důvodu se nová DC často umisťují k jaderným elektrárnám a malým modulárním reaktorům. Ty jim mohou poskytnout stabilní energii, která je zapotřebí kvůli nepřetržitému trénování nových AI modelů.~\cite{chen_electricity_2025, cruzes_data_2025}

Důležitou roli hrají také konektivita a latence. Zde dochází k dělení podle typu zátěže: 

\begin{itemize}
    \item Tradiční cloud a AI inference vyžadují nízkou latenci (okamžitou odezvu uživatelům)~\cite[s.~33]{iea_energy_2025}, a proto se musí stavět poblíž velkých metropolí, nebo přímo v místech hlavní optické sítě~\cite{pilz_compute_2023}, jako je severní Virginie. 
    \item Trénink modelů naopak není citlivý na latenci vůči běžnému uživateli, proto mohou být tato obří centra stavěna mimo hlavní komunikační sítě~\cite{iea_energy_2025}, kde je levnější půda a dostatek energie. Tento trend lze pozorovat na území Texasu a Arizony.
\end{itemize}

\subsection{Legislativní faktory}
\label{subsec:pestle_le}

Legislativní faktory byly z větší části umístěny do jiných oblastí, například způsob získání stavebního povolení spadá do ekonomických faktorů, stejně jako daňové úlevy a jiné iniciativy. Dále by sem mohla být zařazena zmíněná datová suverenita, která je však spíše politickým než legislativním problémem. Z pohledu legislativy je však důležitá otázka, zda a jakým způsobem budou datová centra regulována. V současnosti se již objevují specifické zákony, které se týkají výstavby a provozu datových center. Například v Oregonu byl prosazen zákon, který datová centra zavazuje k alespoň 10letému provozu a zakazuje socializaci nákladů na infrastrukturu (vizte kapitolu 2).~\cite[s.~19]{ginzburg-ganz_technical_2025}

\subsection{Environmentální faktory}
\label{subsec:pestle_en}

Přírodní podmínky mají obrovský dopad na spotřebu energie potřebné k chlazení (PUE). Preferují se lokality s chladnějším a sušším podnebím, které umožňují využívat chlazení venkovním vzduchem (free cooling/economization) po většinu roku~\cite{barroso_data_2026, geng_data_2014}. Neméně kritická je dostupnost vody pro odpařovací chladicí systémy. Stavba datových center vyvolává politické i společenské konflikty, protože centra přímo soutěží o vodní zdroje s místními obyvateli a zemědělstvím~\cite[s.~4]{shah_four_2026}.

Druhým environmentálním faktorem je stabilita a bezpečnost prostředí. Zajištění nepřetržitého chodu znamená vyhnout se oblastem s vysokým rizikem katastrof - lokality se zkoumají z hlediska rizika zemětřesení, záplavových zón, hurikánů, tornád, sesuvů půdy či tsunami.~\cite{tatnall_reflections_2012,geng_data_2014}

\section{Shrnutí}
\label{sec:shrnuti}

%%%%%%%%%%%%%%%%%%%%%%%%%%%%%%%%%%%%%%%%%%%%%%%%%%%%%%%%%%%%%%%%%%%%%%%%%%%%%%%%
% KAPITOLA 6: SHRNUTÍ
% - Stručně opakuje předchozí sekce
% - Je to to nejdůležitější, co je potřeba pro pochopení následujících kapitol
%%%%%%%%%%%%%%%%%%%%%%%%%%%%%%%%%%%%%%%%%%%%%%%%%%%%%%%%%%%%%%%%%%%%%%%%%%%%%%%%

Datová centra prošla za posledních 30 let radikální proměnou. Od sálových počítačů a nezávislých firemních serveroven v 90. letech se průmysl na počátku tisíciletí přesunul k tzv. WSC (Warehouse Scale Computers). Tyto obří komplexy (zpopularizované společnostmi jako Google) propojily statisíce levných CPU serverů do jednoho fungujícího celku, aby obsloužily globální cloudové služby a vyhledávání. Dnešní nástup umělé inteligence však odstartoval novou éru - WSC se mění na AI Factories, což jsou vysoce specializované superpočítače navržené nikoliv pro běh mnoha nezávislých úloh, ale pro masivně paralelní trénink obřích modelů.~\cite{barroso_data_2026,geng_data_2014,barroso_brief_2021}

Tradiční servery postavené na univerzálních procesorech (CPU) přestaly pro AI stačit. Výpočetní zátěž převzaly specializované akcelerátory, jako jsou GPU a TPU, které jsou extrémně efektivní pro maticové operace. S tím však raketově vzrostla spotřeba samotných čipů (TDP) a nutnost umisťovat tyto čipy fyzicky co nejblíže k sobě kvůli zkrácení komunikační latence. To vedlo k extrémní koncentraci výkonu a obrovskému nárůstu hustoty racků (z tradičních 5-10~kW na 30-100+~kW).~\cite{chen_electricity_2025, barroso_data_2026, sheng_power_2026}

Kvůli tomu spotřeba energie datových center dramaticky roste. AI datová centra vstupují do gigawattové éry a zvyšují nároky na metriku PUE (Power Usage Effectiveness). Zásadním problémem pro elektrickou síť ale není jen celkový objem, nýbrž pulzní chování AI zátěže - synchronizované výpočty tisíců GPU generují obrovské výkyvy spotřeby v řádech desítek megawattů během milisekund, což ohrožuje frekvenční stabilitu sítě a zhoršuje kvalitu dodávek (harmonické zkreslení, problikávání). Jedno 200MW datové centrum může ročně spotřebovat přibližně 1,4~TWH elektřiny, což je podobně jako Karlovarský kraj.~\cite{chen_electricity_2025, sheng_power_2026, energeticky_regulacni_urad_spotreba_2024}

Výkonovou hustotu nad 30~kW na rack už nelze spolehlivě uchladit vzduchem. Nastupuje proto kapalinové chlazení (D2C nebo imerzivní), které je nákladnější, ale energeticky mnohem efektivnější. Na druhou stranu vyžaduje obrovské množství vody, která se odpařuje v chladicích věžích. Datová centra navíc nepřímo odebírají vodu při výrobě jimi spotřebované elektřiny v elektrárnách. Průměrné hyperscale centrum může denně spotřebovat až 2 miliony litrů vody, což je podobně jako 6 500 amerických domácností.~\cite{li_making_2025,shah_four_2026,jegham_how_2025,cruzes_data_2025}

Při vybírání lokality pro výstavbu AI datového centra je absolutní prioritou kapacita a propustnost elektrické sítě (často se čeká roky na připojení). Hledají se místa s levnou a stabilní bezuhlíkovou energií (blízkost větrných, solárních či jaderných elektráren), dostatečným zdrojem vody pro chlazení, páteřní optickou sítí a příznivým legislativním i daňovým prostředím.~\cite{lin_exploding_2024,chen_electricity_2025,pilz_compute_2023,geng_data_2014}